\chapter{Resultados}\label{cap:resultados}

\TextoModelo{Apresente as observações e medições de forma organizada, evitando antecipar toda a interpretação. Use a mesma ordem lógica dos métodos ou dos objetivos específicos.}

\section{Unidades geológicas}

\TextoModelo{Descreva os resultados de mapeamento, petrografia, estratigrafia ou caracterização de materiais. Diferencie observação de interpretação.}

\begin{table}[H]
\centering
\caption{Exemplo de síntese das unidades reconhecidas}\label{tab:unidades}
\begin{tabularx}{\textwidth}{l X X}
\toprule
Unidade & Características de campo & Interpretação preliminar \\
\midrule
U1 & \TextoModelo{litologia, textura, estruturas e contatos} & \TextoModelo{ambiente ou processo} \\
U2 & \TextoModelo{litologia, textura, estruturas e contatos} & \TextoModelo{ambiente ou processo} \\
\bottomrule
\end{tabularx}
\FonteTCC{Elaboração própria}
\end{table}

\section{Dados estruturais}

\TextoModelo{Apresente famílias de estruturas, estatísticas descritivas, estereogramas, relações de corte e indicadores cinemáticos.}

\section{Resultados analíticos}

\TextoModelo{Apresente os resultados laboratoriais com unidades, precisão, limites analíticos e identificação inequívoca das amostras.}

\figuraplaceholder{Perfil geológico esquemático da área de estudo}{Elaboração própria}{fig:perfil-geologico}
